% GENERATED FILE. Edit canonical lessons, then run npm run generate:books.
% canonical-source-hash: fnv1a64:59e7719b2a747dd7
% canonical-lessons: FR-C273-veille, FR-C273-lendemain, FR-C273-surlendemain, FR-C273-dernierement, FR-C273-prochainement

\chapter{La veille et le lendemain}
\label{ch:fr-la-veille-et-le-lendemain}
\hlchaptermodality{273}

\begin{chapteropening}
\textbf{By the end of this chapter:} \emph{I can place an event the day before or the day after another, two days later, and say lately or soon.}

Complete the last lesson of chapter 273: I can place an event the day before or the day after another, two days later, and say lately or soon.
\end{chapteropening}

\section[la veille]{la veille — the day before, the eve}

\label{lesson:FR-C273-veille}



\begin{quote}
Before the new one: say the French for until then, then the French for since then.
\end{quote}

\subsection*{You'll want to know: la veille}
\textbf{la veille} — \textquotedblleft{}the day before, the eve\textquotedblright{}.

\begin{quote}
La veille de son examen, elle n'a pas dormi.
\end{quote}

\begin{quote}
\emph{The day before her exam, she did not sleep.}
\end{quote}

\emph{La veille au soir} — the evening before. The verb \emph{veiller} means to stay up.

\begin{cousinweb}
\emph{veille} is inherited from Latin \emph{vigilia} \textquotedblleft{}wakefulness, a night watch\textquotedblright{} (EtymDB), built on \emph{vigil} \textquotedblleft{}awake\textquotedblright{}: the night of waiting before a feast.
\end{cousinweb}

\begin{grammarlens}[title={What you've built}]
One more word for telling, explaining and describing at B1.
\end{grammarlens}

\subsection*{Your turn}
\begin{itemize}
\raggedright
  \item \emph{Say it:} \emph{la veille}
  \item \emph{Say it:} \emph{la veille}, once more
  \item \emph{Say it:} a sentence of your own with \emph{la veille}
  \item \emph{Recall:} say the French for until then, then the French for since then, then say \emph{la veille} again
\end{itemize}

\begin{tcolorbox}[breakable,colback=teal!4,colframe=teal!35!black,title={Before you move on}]
What does la veille mean? (\textquotedblleft{}The day before, the eve\textquotedblright{}.) Say it once more.
\end{tcolorbox}
\section[le lendemain]{le lendemain — the next day, the day after}

\label{lesson:FR-C273-lendemain}



\begin{quote}
Before the new one: say the French for since then, then the French for the day before.
\end{quote}

\subsection*{You'll want to know: le lendemain}
\textbf{le lendemain} — \textquotedblleft{}the next day, the day after\textquotedblright{}.

\begin{quote}
Le lendemain, il a téléphoné pour s'excuser.
\end{quote}

\begin{quote}
\emph{The next day, he phoned to apologise.}
\end{quote}

It looks forward from a day in the past; from today, say \emph{demain}.

\begin{cousinweb}
\emph{lendemain} grew out of \emph{l'endemain}, built on \emph{demain}, which is inherited from Latin \emph{dē māne} \textquotedblleft{}from the morning\textquotedblright{} (Lewis \& Short).
\end{cousinweb}

\begin{grammarlens}[title={What you've built}]
One more word for telling, explaining and describing at B1.
\end{grammarlens}

\subsection*{Your turn}
\begin{itemize}
\raggedright
  \item \emph{Say it:} \emph{le lendemain}
  \item \emph{Say it:} \emph{le lendemain}, once more
  \item \emph{Say it:} a sentence of your own with \emph{le lendemain}
  \item \emph{Recall:} say the French for since then, then the French for the day before, then say \emph{le lendemain} again
\end{itemize}

\begin{tcolorbox}[breakable,colback=teal!4,colframe=teal!35!black,title={Before you move on}]
What does le lendemain mean? (\textquotedblleft{}The next day, the day after\textquotedblright{}.) Say it once more.
\end{tcolorbox}
\section[le surlendemain]{le surlendemain — two days later, the day after next}

\label{lesson:FR-C273-surlendemain}



\begin{quote}
Before the new one: say the French for the day before, then the French for the next day.
\end{quote}

\subsection*{You'll want to know: le surlendemain}
\textbf{le surlendemain} — \textquotedblleft{}two days later, the day after next\textquotedblright{}.

\begin{quote}
Il est arrivé le lundi et il est reparti le surlendemain.
\end{quote}

\begin{quote}
\emph{He arrived on the Monday and left again two days later.}
\end{quote}

\emph{Sur} + \emph{lendemain}: one day beyond the next day.

\begin{cousinweb}
\emph{surlendemain} is built on \emph{lendemain} with \emph{sur} \textquotedblleft{}over, beyond\textquotedblright{}; \emph{demain} is inherited from Latin \emph{dē māne} \textquotedblleft{}from the morning\textquotedblright{} (Lewis \& Short).
\end{cousinweb}

\begin{grammarlens}[title={What you've built}]
One more word for telling, explaining and describing at B1.
\end{grammarlens}

\subsection*{Your turn}
\begin{itemize}
\raggedright
  \item \emph{Say it:} \emph{le surlendemain}
  \item \emph{Say it:} \emph{le surlendemain}, once more
  \item \emph{Say it:} a sentence of your own with \emph{le surlendemain}
  \item \emph{Recall:} say the French for the day before, then the French for the next day, then say \emph{le surlendemain} again
\end{itemize}

\begin{tcolorbox}[breakable,colback=teal!4,colframe=teal!35!black,title={Before you move on}]
What does le surlendemain mean? (\textquotedblleft{}Two days later, the day after next\textquotedblright{}.) Say it once more.
\end{tcolorbox}
\section[dernièrement]{dernièrement — lately, recently}

\label{lesson:FR-C273-dernierement}



\begin{quote}
Before the new one: say the French for the next day, then the French for two days later.
\end{quote}

\subsection*{You'll want to know: dernièrement}
\textbf{dernièrement} — \textquotedblleft{}lately, recently\textquotedblright{}.

\begin{quote}
Dernièrement, je dors mal.
\end{quote}

\begin{quote}
\emph{Lately, I have been sleeping badly.}
\end{quote}

Close to \emph{récemment}; it goes with the present or the passé composé.

\begin{cousinweb}
\emph{dernièrement} is built on \emph{dernier} \textquotedblleft{}last\textquotedblright{}, from Old French \emph{derrenier}, usually traced to Latin \emph{dē retrō} \textquotedblleft{}from behind\textquotedblright{} (Lewis \& Short).
\end{cousinweb}

\begin{grammarlens}[title={What you've built}]
One more word for telling, explaining and describing at B1.
\end{grammarlens}

\subsection*{Your turn}
\begin{itemize}
\raggedright
  \item \emph{Say it:} \emph{dernièrement}
  \item \emph{Say it:} \emph{dernièrement}, once more
  \item \emph{Say it:} a sentence of your own with \emph{dernièrement}
  \item \emph{Recall:} say the French for the next day, then the French for two days later, then say \emph{dernièrement} again
\end{itemize}

\begin{tcolorbox}[breakable,colback=teal!4,colframe=teal!35!black,title={Before you move on}]
What does dernièrement mean? (\textquotedblleft{}Lately, recently\textquotedblright{}.) Say it once more.
\end{tcolorbox}
\section[prochainement]{prochainement — soon, shortly}

\label{lesson:FR-C273-prochainement}



\begin{quote}
Before the new one: say the French for two days later, then the French for lately.
\end{quote}

\subsection*{You'll want to know: prochainement}
\textbf{prochainement} — \textquotedblleft{}soon, shortly\textquotedblright{}.

\begin{quote}
Le magasin va ouvrir prochainement.
\end{quote}

\begin{quote}
\emph{The shop is going to open soon.}
\end{quote}

A word of announcements and posters: \emph{prochainement au cinéma}.

\begin{cousinweb}
\emph{prochainement} is built on \emph{prochain} \textquotedblleft{}next\textquotedblright{}, inherited from Vulgar Latin \emph{propeānus}, built on Latin \emph{prope} \textquotedblleft{}near\textquotedblright{} (Lewis \& Short).
\end{cousinweb}

\begin{grammarlens}[title={What you've built}]
One more word for telling, explaining and describing at B1.
\end{grammarlens}

\subsection*{Your turn}
\begin{itemize}
\raggedright
  \item \emph{Say it:} \emph{prochainement}
  \item \emph{Say it:} \emph{prochainement}, once more
  \item \emph{Say it:} a sentence of your own with \emph{prochainement}
  \item \emph{Recall:} say the French for two days later, then the French for lately, then say \emph{prochainement} again
\end{itemize}

\begin{tcolorbox}[breakable,colback=teal!4,colframe=teal!35!black,title={Before you move on}]
What does prochainement mean? (\textquotedblleft{}Soon, shortly\textquotedblright{}.) Say it once more.
\end{tcolorbox}
